\section{Threat Model}
\label{sec:threat}

\subsection{Adversary Capabilities}

We adopt a Kerckhoffs-compliant threat model~\cite{kerckhoffs1883}: the adversary possesses
\emph{full knowledge} of the \sentinel{} defense architecture, including all
detection patterns, layer mechanisms, escalation rules, and classification
taxonomies. Security rests not on architectural secrecy but on the
confidentiality of four classes of ephemeral state:

\begin{enumerate}
    \item \textbf{Ephemeral cryptographic keys} rotated per session and used
          for inter-layer attestation.
    \item \textbf{Active canary probes} injected by the \tsa{} layer to
          measure downstream integrity.
    \item \textbf{Activation baselines} computed over sliding windows and
          unique to each deployment instance.
    \item \textbf{Negative-selection detector sets} generated stochastically
          and never fully observable by an external party.
\end{enumerate}

\noindent
This separation ensures that publishing the full architectural specification
(as we do in this paper) does not weaken operational deployments.

\subsection{Attack Categories}

We organize adversarial inputs into fifteen categories spanning the known
attack surface of large language model deployments.
Table~\ref{tab:attack-categories} summarizes each category and the number of
simulated attack prompts used in our evaluation (Section~\ref{sec:evaluation}).

\begin{table}[t]
\centering
\caption{Attack categories and simulation counts used for \sentinel{} evaluation.
         Total corpus: 250{,}000 adversarial prompts.}
\label{tab:attack-categories}
\small
\begin{tabular}{@{}lr@{}}
\toprule
\textbf{Category} & \textbf{Simulations} \\
\midrule
Direct Injection              & 25{,}000 \\
Indirect Injection            & 25{,}000 \\
Multi-turn Crescendo~\cite{russinovich2024} & 20{,}000 \\
Encoding / Obfuscation        & 20{,}000 \\
Role-play / Persona           & 20{,}000 \\
Tool Abuse / Agentic          & 20{,}000 \\
Data Exfiltration             & 15{,}000 \\
Social Engineering            & 15{,}000 \\
Semantic Equivalence          & 15{,}000 \\
Steganographic                & 12{,}000 \\
Model-Level Compromise~\cite{hubinger2024} & 10{,}000 \\
Cross-boundary Trust          & 10{,}000 \\
Novel / Zero-day              & 13{,}000 \\
Multi-modal                   & 10{,}000 \\
Adversarial ML                & 10{,}000 \\
\midrule
\textbf{Total}                & \textbf{250{,}000} \\
\bottomrule
\end{tabular}
\end{table}

\paragraph{Category design rationale.}
Categories were selected to cover both well-studied vectors (direct and
indirect injection) and emerging threats (steganographic channels,
model-level compromise via sleeper agents~\cite{hubinger2024}).  The
\emph{Novel / Zero-day} category contains prompts generated by an
unconstrained red-team model with no access to \sentinel{} internals,
providing an estimate of robustness against previously unseen attack
patterns.

\subsection{Mutation Strategy}

A realistic adversary does not submit a single static payload; instead, they
iteratively mutate failed attempts.  To model this adaptive behavior, we
define five mutation operators and measure the \pasr{} degradation each
induces on \sentinel{} when applied to initially-detected prompts.
Table~\ref{tab:mutations} reports results averaged across all fifteen attack
categories.

\begin{table}[t]
\centering
\caption{Mutation operators and their effect on \sentinel{} detection rate
         (\pasr{} degradation relative to unmutated baseline).}
\label{tab:mutations}
\small
\begin{tabular}{@{}llr@{}}
\toprule
\textbf{Mutation Type} & \textbf{Description} & \textbf{Degradation} \\
\midrule
Lexical      & Synonym substitution        & $-8.7\%$  \\
Structural   & Reorder clauses / sentences & $-6.1\%$  \\
Encoding     & Switch encoding schemes     & $-14.5\%$ \\
Context      & Change cover story          & $-12.3\%$ \\
Hybrid       & Combine $\geq 2$ operators  & $-18.2\%$ \\
\bottomrule
\end{tabular}
\end{table}

\noindent
Hybrid mutations pose the greatest challenge, yet even under worst-case
composition the \mire{} layer's semantic analysis and the \tsa{} layer's
behavioral monitoring jointly maintain detection rates above the operational
threshold (Section~\ref{sec:evaluation}).

\subsection{Impossibility Results}

Two fundamental impossibility results bound what any defense---including
\sentinel{}---can achieve.

\paragraph{Indistinguishability of backdoored models.}
Goldwasser and Kim~\cite{goldwasser2022} show that no polynomial-time
algorithm can reliably distinguish a backdoored neural network from a clean
one, given only black-box query access.  This result implies that purely
input--output testing is insufficient for detecting model-level compromise;
\sentinel{} therefore supplements query-time analysis with activation-space
monitoring at the \tsa{} layer, operating on internal representations rather
than solely on outputs.

\paragraph{Semantic identity limitation.}
Informally, for any classifier~$C$ over natural-language requests, there
exist semantically equivalent request pairs $(r, r')$ such that $C$ must
misclassify at least one member of the pair---either producing a false
positive on a benign request or a false negative on a malicious one.  This
follows from the undecidability of semantic equivalence in unrestricted
natural language and motivates the multi-layer design of \sentinel{}: no
single classification boundary suffices, so the lattice composes
complementary detection surfaces (\pasr{}, \mire{}, \tsa{}) whose error
regions are empirically near-disjoint (Section~\ref{sec:evaluation}).

\paragraph{Operational stance.}
\sentinel{} does not claim to circumvent these theoretical limits.  Rather,
the architecture is designed to \emph{operate within them}: layered
redundancy ensures that the residual attack surface left by any individual
impossibility result is covered by an orthogonal detection mechanism,
yielding aggregate detection rates that approach---but cannot
reach---perfection.
